\chapter{Akbarian Metaphysics and the Self-Disclosure of Wujud}
\label{ch:akbarian}

\epigraph{\itshape The cosmos is His self-disclosure to Himself in
the loci of His self-disclosure.}{--- Ibn al-`Arabi, paraphrased
from the \emph{Futuhat al-Makkiyyah}}

\noindent
This chapter develops the Akbarian metaphysical tradition --- the
philosophical school descending from Muhyi al-Din ibn al-`Arabi
(d.~638/1240) and continued through Mulla Sadra (d.~1641) and Shah
Wali Allah Dehlawi (d.~1762). The chapter is the philosophical
climax of Part II. The four preceding chapters have developed
contemporary intellectual resources --- quantum interpretations,
ergodicity economics, analytical idealism, non-computability ---
that approach, from outside, the metaphysical position that the
Akbarian tradition has articulated from inside for nearly eight
hundred years.

The chapter is the longest in Part II because the Akbarian system
is genuinely complex. It is also philosophically the richest,
combining ontological precision with religious depth in a
combination that has few parallels in the Western tradition. The
chapter is the conceptual bridge between the modern analytical
apparatus we have built up and the worked treatments of Part III
that will use the apparatus to articulate specific Islamic economic
claims.

The chapter assumes basic familiarity with the philosophical issues
of \xref{ch:analytical-idealism} but no prior knowledge of Islamic
philosophy. The exposition proceeds slowly enough that the reader
new to the tradition can follow; readers familiar with the tradition
will find the treatment compressed but, the author hopes,
illuminating.

The chapter proceeds in twelve sections.
\xref{sec:ibn-arabi-life} introduces Ibn al-`Arabi and his place in
the tradition.
\xref{sec:wujud-foundational} develops the central concept of
\arb{wujud} (being).
\xref{sec:wahdat-al-wujud-detail} treats the doctrine of
\arb{wahdat al-wujud} (the unity of being), distinguishing it
carefully from pantheism.
\xref{sec:divine-names} treats the divine names as principles of
articulation.
\xref{sec:tajalli-mechanism} develops \arb{tajalli} (self-disclosure)
as the active principle of creation.
\xref{sec:imaginal-world} treats the imaginal world (\arb{`alam
al-mithal}) as an ontological domain.
\xref{sec:perfect-human} treats the perfect human (\arb{insan
al-kamil}).
\xref{sec:mulla-sadra} treats Mulla Sadra's substantial motion.
\xref{sec:shah-wali-allah} treats Shah Wali Allah's economic
theology.
\xref{sec:contemporary-correspondence} draws the correspondence with
the contemporary frameworks.
\xref{sec:akbarian-implications} previews the implications for the
eight worked treatments of Part III.

\section{Ibn al-`Arabi: life and tradition}
\label{sec:ibn-arabi-life}

\subsection{The life}

Muhyi al-Din ibn al-`Arabi was born in Murcia, in the Muslim
Andalusia of the Almohad period, in 560 AH (1165 CE). His family
was distinguished; his early education was extensive. By his
adolescence he had received what he called his ``opening''
(\arb{fath}): a series of mystical experiences that shaped his
subsequent intellectual life. He traveled to Tunis, Cairo, Mecca,
Konya, and finally Damascus, where he settled and wrote much of his
mature work. He died in 638 AH (1240 CE) and is buried in Damascus.

His written corpus is enormous. The \textit{Futuhat al-Makkiyyah}
(Meccan Revelations), his magnum opus, runs to four large volumes
in modern editions, comprising 560 chapters that cover essentially
every topic in Islamic intellectual life. The \textit{Fusus
al-Hikam} (Bezels of Wisdom), much shorter and more concentrated,
treats the metaphysical content most systematically. Other works
include the \textit{Tarjuman al-Ashwaq} (Interpreter of Desires,
poetry with prose commentary), the \textit{Risalah al-Anwar}
(Treatise on Lights), and dozens of treatises on specific topics.

The corpus has been the subject of continuous scholarly engagement
since Ibn `Arabi's death. The major Akbarian commentators include
Sadr al-Din al-Qunawi (d.~1274), Ibn `Arabi's stepson and direct
student; Mu'ayyad al-Din al-Jandi (d.~1300); `Abd al-Razzaq
al-Kashani (d.~1335); Da'ud al-Qaysari (d.~1350); and many others
through the centuries. The tradition spread from Andalusia and the
Eastern Mediterranean across the Muslim world; by the seventeenth
century the Akbarian framework was the dominant metaphysical
framework in Mughal India, Safavid Iran, and the Ottoman Empire.

\subsection{The school}

The Akbarian school is named for Ibn `Arabi's title \arb{al-Shaykh
al-Akbar} (the Greatest Master). The school is not a sect or a
religious order; it is a philosophical and contemplative tradition
that includes thinkers from various Sufi orders, from various
schools of fiqh, and from both the Sunni and Shi`i communities.
Ibn `Arabi himself was a Sunni Muslim of (broadly) the Maliki
school; his philosophical influence has spread far beyond his own
denominational identity.

The school's central commitments are the doctrines we treat in this
chapter: \arb{wahdat al-wujud}, the divine names, \arb{tajalli},
the imaginal world, and the perfect human. These commitments are
articulated with great philosophical care; they are not merely
mystical assertions but worked-out philosophical positions that
respond to specific philosophical problems.

The contemporary academic engagement with the Akbarian tradition,
in Western languages, is comparatively recent. Henry Corbin's work
in French (mid-twentieth century) was foundational. The most
extensive English-language scholarship is by William Chittick
(SUNY Stony Brook), particularly his \textit{The Sufi Path of
Knowledge} (1989) and \textit{The Self-Disclosure of God} (1998).
Toshihiko Izutsu's \textit{Sufism and Taoism} (1984) is a
comparative classic. Sachiko Murata's \textit{The Tao of Islam}
(1992) treats the tradition's account of complementarity. The
Akbarian tradition is, in the contemporary academic context, a
philosophical resource that is finally being made widely available
in rigorous translation and exposition.

\subsection{The location of the tradition in our project}

The Akbarian tradition has a specific role in this book. It is the
classical Islamic articulation of the ontology that the contemporary
frameworks of \xref{ch:analytical-idealism} and
\xref{ch:non-computability} approach from outside. Where Kastrup's
analytical idealism is a contemporary defense of the
mind-fundamentality position, Ibn `Arabi's metaphysics is the same
position articulated within the Islamic tradition with a different
vocabulary, different philosophical antecedents, and (importantly)
a richer theological content.

The chapter therefore treats the Akbarian tradition seriously as
philosophy, with full attention to its conceptual content, while
also acknowledging the theological content that it adds beyond the
contemporary philosophical framework. The reader who is not Muslim
is invited to engage with the philosophical content; the reader who
is Muslim will recognize the tradition as their own and find,
perhaps, a contemporary articulation that connects it to current
philosophical work.

\section{Wujud: the foundational ontological category}
\label{sec:wujud-foundational}

\subsection{What wujud is}

\begin{definition}[\arb{Wujud}]
\label{def:wujud}
The Arabic term \arb{wujud} (root \textit{w-j-d}, ``to find,''
``to be present'') is variously translated as ``being,''
``existence,'' ``finding,'' ``presence.'' In Akbarian metaphysics,
\arb{wujud} is the foundational ontological category: it names
that-which-is, prior to and underlying any particular determinate
existence. \arb{Wujud} is not a property; it is the ground in which
all properties subsist.
\end{definition}

The English ``existence'' is misleading as a translation. The
modern English ``existence'' is typically a binary property: a
thing either exists or does not. \arb{Wujud}, in the Akbarian
sense, is not like this. \arb{Wujud} is a single reality with
internal structure; what we ordinarily call ``things'' are modes,
gradations, or loci of \arb{wujud}.

The classical analogy in the tradition is light. Light is one
reality; it manifests in different colors, different intensities,
different forms; the colors and forms are not separate things but
modes of the single light. \arb{Wujud} is the same: one reality, with
the apparent multiplicity of things being modes of it.

\subsection{The grades of wujud}

The Akbarian tradition distinguishes several grades or levels of
\arb{wujud}. The standard enumeration includes:

\begin{itemize}[leftmargin=2em]
    \item \emph{Absolute \arb{wujud}}: \arb{wujud} as such, without
    qualification; the divine being, identified with Allah.
    \item \emph{Universal \arb{wujud}}: the manifestation of
    absolute \arb{wujud} in the universal forms (the divine names
    and attributes).
    \item \emph{The world of intelligences}: the realm of pure
    intellect, the angelic order.
    \item \emph{The world of souls}: the realm of psychic reality,
    including human souls.
    \item \emph{The imaginal world (\arb{`alam al-mithal})}: the
    intermediate ontological domain; treated in \xref{sec:imaginal-world}.
    \item \emph{The world of bodies}: the material world, the realm
    of physical things.
\end{itemize}

The grades are not discontinuous levels. They are gradations of a
single \arb{wujud}, with the lower grades being attenuations or
manifestations of the higher. The world of bodies is the most
attenuated grade; absolute \arb{wujud} is the source from which
all the others derive.

\subsection{The contemporary correspondence}

The Akbarian conception of \arb{wujud} as a single reality with
internal structure corresponds, structurally, to Kastrup's
Mind-at-Large of \xref{ch:analytical-idealism}. Both are single
fundamental ontological realities; both have internal structure
that gives rise to the apparent multiplicity of things; both are
presupposed by, rather than constructed from, the things that
appear within them.

The main difference, which we develop throughout the chapter, is
that the Akbarian framework has more articulated structure. The
divine names provide a worked-out account of \emph{how} the
underlying unity differentiates into determinate things. The
analytical-idealist framework, in its current form, has the unity
and the dissociated alters but does not articulate the
intermediate structure that produces specific kinds of beings.

\section{Wahdat al-wujud: the unity of being}
\label{sec:wahdat-al-wujud-detail}

\subsection{The doctrine}

\begin{conceptbox}[\arb{Wahdat al-wujud}]
\label{conc:wahdat}
\arb{Wahdat al-wujud} --- the unity of being --- is the central
metaphysical doctrine of the Akbarian tradition. It asserts that
there is, ultimately, only one \arb{wujud}, and that all created
things are not independent existents but loci of divine
self-disclosure. Created things are real, but their reality is
borrowed; their being is not their own.
\end{conceptbox}

The doctrine is the Akbarian articulation of the Quranic principle
\arb{la ilaha illa Allah} --- ``there is no being but Allah.''
Whereas the standard theological reading takes the principle to
assert the exclusivity of worship (no being is worthy of worship
but Allah), the Akbarian reading takes it to assert the exclusivity
of being (no \emph{wujud} but the divine \emph{wujud}). Created
things have being only in the derivative sense that they manifest
the divine being; they have no independent ontological status.

\subsection{Distinction from pantheism}

The most common misreading of \arb{wahdat al-wujud}, both in the
Islamic and the Western philosophical literature, is to identify it
with pantheism. This is a serious error, and the Akbarian tradition
has consistently corrected it.

\begin{conceptbox}[Pantheism vs.\ \arb{wahdat al-wujud}]
\textbf{Pantheism}: ``Everything is God.'' The world, in its full
multiplicity of things, is the totality of divine being. God is the
sum of the things; the divine is exhausted by the world.

\textbf{Wahdat al-wujud}: ``There is no being but God.'' The world,
in its multiplicity, is the appearance of divine being in finite,
conditioned modes. The world is real \emph{as appearance} and
unreal \emph{as independent existent}. God is not the sum of the
things; God is the sole reality of which the things are
manifestations.
\end{conceptbox}

The difference is fundamental. Pantheism collapses God into the
world (and is heretical from the Islamic perspective). \arb{Wahdat
al-wujud} maintains a clear distinction: God is the absolute
reality; the world is the appearance of that reality in
gradated, conditioned modes. The world is therefore not God in any
straightforward sense; the world is the self-disclosure of God to
Himself.

The closest Western analogue is not Spinoza's pantheism but
Berkeley's idealism (or, in contemporary form, Kastrup's analytical
idealism). For Berkeley, the world consists of ideas in minds (most
fundamentally, in the divine mind); the world is real as appearance
to the divine mind, not as independent matter. This is structurally
much closer to \arb{wahdat al-wujud} than pantheism is.

\subsection{The strength of the claim}

It is worth being explicit about how strong the Akbarian claim is.
The doctrine asserts that:

\begin{itemize}[leftmargin=2em]
    \item There is only one being. The apparent existence of many
    beings is the appearance of the one being in many modes.
    \item This being is not equal to or the sum of the apparent
    things. It is their source and their substrate.
    \item Created things have no being of their own. Their reality is
    entirely borrowed from the one being.
    \item The relationship between the one being and its modes is
    not causation in the standard sense. The modes are not produced
    \emph{by} the one being as a separate effect; they are the one
    being \emph{appearing as} those modes.
\end{itemize}

This is a substantial metaphysical position. The Akbarian tradition
defends it with extensive philosophical argumentation, much of it
preserved in the commentary literature. The arguments include:

\begin{itemize}[leftmargin=2em]
    \item The argument from divine simplicity: if God is absolute and
    unconditioned, there cannot be a real ontological gap between
    God and other beings, since such a gap would constitute a
    determination of God by something outside Him.
    \item The argument from the meaning of being: the term
    ``being'' applies univocally to all things only if there is a
    single underlying reality of which they are modes; otherwise it
    is equivocal.
    \item The argument from the structure of self-knowledge: divine
    self-knowledge is exhaustive; the things that exist are divine
    self-knowledge actualized.
\end{itemize}

The arguments are technical and we do not rehearse them in detail.
The point is that \arb{wahdat al-wujud} is a worked-out
philosophical position, not a mystical assertion. The Akbarian
philosophers defended it with the same rigor that contemporary
analytic philosophers bring to debates about idealism.

\section{The divine names: principles of articulation}
\label{sec:divine-names}

\subsection{The Quranic basis}

The Quran lists, in various passages, the most beautiful names of
God (\arb{al-asma' al-husna}). The classical enumeration gives 99
names, although the actual count varies by exegetical tradition.
Examples: \arb{ar-Rahman} (the Merciful), \arb{ar-Razzaq} (the
Provider), \arb{al-Khaliq} (the Creator), \arb{al-`Alim} (the
Knowing), \arb{al-Hakim} (the Wise), \arb{al-Latif} (the Subtle),
\arb{al-Qadir} (the Powerful), and so on.

The standard theological treatment takes these as descriptors of
divine attributes. God is merciful (so the name \arb{ar-Rahman} is
true of Him); God is the provider (so \arb{ar-Razzaq} is true);
and so on. The names are predicates of the divine subject.

The Akbarian treatment is much richer. The names are not just
predicates; they are the structural principles by which absolute
\arb{wujud} unfolds into determinate being. Each name is a specific
mode of divine self-disclosure; together, the names constitute the
articulated structure of created reality.

\subsection{Names as modes of divine self-disclosure}

\begin{conceptbox}[The divine names as modes]
\label{conc:divine-names}
The divine names are not merely descriptors. They are the structural
principles by which absolute \arb{wujud} unfolds into the
determinate beings of creation. Each name is a specific mode of
divine self-disclosure; the totality of divine names is the
articulated structure of created reality.
\end{conceptbox}

The conception is foreign to most modern theology and requires
careful unpacking. To say that \arb{ar-Razzaq} (the Provider) is a
divine name is, on the standard theology, to say that God provides.
On the Akbarian conception, it is to say something more: that
``providing'' is one of the structural principles of reality, that
this principle has a specific character, and that it produces
specific kinds of beings (those that receive provision) when it
manifests.

The grain of wheat, on the Akbarian conception, is a locus of
several divine names. It manifests \arb{al-Khaliq} (creating
existence). It manifests \arb{ar-Razzaq} (the divine providing
nature, here as nutriment). It manifests \arb{al-Latif} (the
subtle, in the delicacy of the seed). The grain of wheat is what
this specific configuration of divine names looks like when it has
been instantiated as a particular being.

The human being, similarly, is a locus of divine names --- but a
much more comprehensive one. The human can manifest, in principle,
all the divine names. Different individuals manifest them in
different proportions; the perfect human (\arb{insan al-kamil},
treated in \xref{sec:perfect-human}) is the locus that manifests
all of them in equilibrium.

\subsection{The structural significance}

The doctrine of the divine names gives the Akbarian framework a
specific structural advantage over the contemporary
analytical-idealist framework. Where Kastrup's analytical idealism
has the basic ontology (Mind-at-Large, dissociated alters) but does
not articulate the intermediate structure that produces specific
kinds of beings, the Akbarian framework articulates this structure
in detail.

The divine names are the principles by which the unified underlying
\arb{wujud} differentiates into specific kinds. The world is not a
random distribution of dissociated alters; it has structure, and
the structure follows the structure of the divine names. The kinds
of beings that exist (animal, vegetable, mineral, human, angelic)
are not arbitrary; they correspond to the structural possibilities
of divine self-disclosure.

This is a substantial elaboration of the basic idealist position,
and it is what gives the Akbarian framework its capacity to do
specific economic work. The economic claims of Islamic tradition
about provision (\arb{rizq}), about blessing (\arb{barakah}), about
intention (\arb{niyyah}) are all readable as claims about the
manifestation of specific divine names in specific configurations.
We develop this in \xref{sec:akbarian-implications} and in Part III.

\section{Tajalli: the active principle of creation}
\label{sec:tajalli-mechanism}

\subsection{The doctrine}

\begin{definition}[\arb{Tajalli}]
\label{def:tajalli}
\arb{Tajalli} (root \textit{j-l-w}, ``to make manifest,'' ``to
shine forth''), translated as ``theophany'' or ``self-disclosure,''
is the active principle by which absolute \arb{wujud} manifests in
finite loci of existence. Each moment of created existence is a
fresh \arb{tajalli}; the universe is not a created clockwork that
runs on its own but a continuous unfolding of divine
self-disclosure.
\end{definition}

\arb{Tajalli} is the dynamic counterpart of the structural concepts
we have developed. \arb{Wahdat al-wujud} is the structural claim
that there is one being; \arb{tajalli} is the dynamic claim that
this one being is, at every moment, actively manifesting in finite
loci. Without \arb{tajalli}, the world would not exist; the world
exists because the divine self-disclosure is continuously occurring.

\subsection{Continuous creation}

The Akbarian conception of creation is therefore continuous, not
once-and-for-all. The standard Aristotelian conception of God ---
the unmoved mover who set the world in motion and now sustains it
from outside --- is foreign to Akbarian metaphysics. On the Akbarian
conception, God is not setting the world in motion from outside;
God is, at every moment, instantiating the world by self-disclosure.
If the disclosure ceased, the world would not just stop but would
cease to be.

This conception has several theological motivations. It is supported
by Quranic verses such as 55:29: ``Every day He is upon some
matter,'' which assert ongoing divine activity. It is supported by
the hadith literature on divine attention to particulars. It is
supported by the philosophical argument that contingent beings
cannot sustain themselves in being without continuous infusion of
\arb{wujud} from the source.

The conception has also philosophical advantages. It avoids the
deistic problem of a clockwork universe with an absentee God. It
gives a coherent account of why the world has the regularities it
does (the regularities are the structural patterns of divine
self-disclosure). It permits a richer account of divine
interaction with the world (every moment is a divine act, not just
the moments of explicit miracle).

\subsection{Tajalli and contemporary physics}

The Akbarian conception of \arb{tajalli} as continuous self-disclosure
has interesting resonances with contemporary physics, particularly
quantum mechanics. The wave function is not a once-given specification
of reality; it is continuously evolved by the Schr\"odinger
equation. Measurement events --- moments of definite outcome ---
occur continuously throughout the cosmos as quantum systems
interact. Each interaction is, in some interpretations, a fresh
``actualization'' of the underlying possibilities.

The structural parallel is striking. The Akbarian \arb{tajalli} is
the divine self-disclosure that produces the world at each moment;
the quantum measurement event is the actualization of the wave
function at each moment of interaction. Both are dynamic; both
occur continuously; both are required for the world to be as it
appears.

The parallel is structural, not strict identity. We are not claiming
that quantum measurement \emph{is} \arb{tajalli}. We are claiming
that the structure of ``continuous instantiation by an underlying
process'' is articulated in both frameworks, and that the Akbarian
framework provides a metaphysical account of why physical reality
has the dynamic character that the quantum formalism describes.

\subsection{The economic implications}

The continuous-creation conception has direct economic implications.
Wealth, on this view, is not a substance we accumulate. It is a
manifestation that we are, at this moment, receiving. Tomorrow's
wealth will be a fresh manifestation, not a continuation of
today's; the continuity is in the underlying source, not in any
material substance that persists.

This re-frames every economic claim. The Quranic statement that
``upon Allah is the provision of every creature'' (\Quran{11:6})
is, on the Akbarian reading, the assertion that provision is a
manifestation of the divine name \arb{ar-Razzaq} in the locus of
the soul, occurring continuously. The believer's provision is not
generated by their efforts (though their efforts are part of the
manifestation); it is supplied by the underlying source.

To claim absolute possession of wealth (``my wealth'') is, in
Akbarian terms, a metaphysical error: the wealth is not a substance
that persists in my possession; it is a manifestation that I am
receiving. The discipline of zakat, the prohibition of riba, and
the practice of charity are not arbitrary moral rules; they are
acknowledgments of the actual ontological structure of wealth ---
that it is a borrowed reality, continuously received from the source.

We develop the specific economic implications in Part III.

\section{The imaginal world}
\label{sec:imaginal-world}

\subsection{The doctrine}

The Akbarian tradition introduces a distinctive ontological category:
the imaginal world (\arb{`alam al-mithal}), an intermediate
ontological domain between pure spirit (the realm of intelligences
and souls) and pure matter (the world of bodies).

\begin{definition}[\arb{`Alam al-mithal}]
\label{def:imaginal-world}
The \emph{imaginal world} (\arb{`alam al-mithal}) is an intermediate
ontological domain between the spiritual and the material. It is
the realm of forms-with-images: realities that have the structural
character of physical things (extension, color, sound) but are not
themselves physical, and that have the structural character of
spiritual realities (consciousness, value, intentionality) but are
not themselves purely spiritual. Visions, dreams, prophetic
experiences, and certain visual aspects of contemplative practice
are encounters with the imaginal world.
\end{definition}

The doctrine, articulated by Ibn `Arabi and developed by his
commentators, is one of the most distinctive Akbarian contributions.
The imaginal world is not a mere subjective faculty (the
``imagination'' in the modern psychological sense); it is an
ontological domain with its own reality, accessed by the
imagination as a real cognitive faculty.

\subsection{Why the imaginal world is needed}

The doctrine addresses a classical problem: how can immaterial
soul interact with material body? The classical answers in Western
philosophy --- Cartesian interactionism, Spinozistic parallelism,
Leibnizian pre-established harmony --- have generally been
unsatisfying. The interaction problem persists.

The Akbarian answer is the imaginal world. The bridge between
spirit and matter is not a direct interaction across an
ontological gap; it is a graded ontological structure with the
imaginal world as the intermediate domain. Spirit acts on the
imaginal; the imaginal interfaces with matter; matter responds.
The imaginal world is the medium through which spirit and matter
interact.

This account has several attractive features. It is consistent with
the phenomenology of contemplative experience (which routinely
encounters realities that are neither purely spiritual nor purely
material). It is consistent with the tradition's high estimation
of dreams and visions (which are encounters with the imaginal). It
gives a structural account of how the body can be ``ensouled''
without the soul being a separate substance trapped in the body.

\subsection{The imaginal in contemporary philosophy of mind}

The doctrine of the imaginal world has not been seriously engaged
with in contemporary Western philosophy. The Western tradition has
typically treated imagination as either psychological (a faculty of
mind) or fictional (the production of non-existent objects). The
Akbarian conception, that imagination accesses a real ontological
domain, is foreign.

This is one place where the Akbarian framework substantially exceeds
what the contemporary analytical-idealist framework has articulated.
Kastrup's analytical idealism gives us Mind-at-Large and dissociated
alters; it does not have a worked-out account of the intermediate
ontological structure between mental processes and physical
appearance. The Akbarian \arb{`alam al-mithal} provides this
intermediate structure.

For our economic project, the doctrine has specific applications.
The valuation of goods, the meaning of intention, the structure of
trust --- these are all phenomena that have aspects that are
neither purely material (a coin is not just a metal disk) nor purely
mental (the value is not just a belief). They are imaginal in the
Akbarian sense: they have the structural character of both, mediated
by the intermediate ontological domain.

We do not develop this in detail here; it is one of the directions
in which the Akbarian framework can be extended in subsequent
volumes. For the present project, we note the doctrine and its
distinctive contribution.

\section{The perfect human}
\label{sec:perfect-human}

\subsection{The doctrine}

\begin{definition}[\arb{Insan al-kamil}]
\label{def:perfect-human}
The \emph{perfect human} (\arb{insan al-kamil}) is the locus that
manifests all the divine names in equilibrium. The perfect human is
not a divine being; the perfect human is a creature whose locus has
the maximum capacity for divine self-disclosure. In Akbarian
thought, the Prophet Muhammad is the exemplar of \arb{insan al-kamil}.
\end{definition}

The doctrine of \arb{insan al-kamil} is the Akbarian articulation of
the special place of the human being in creation. Other creatures
are loci of specific configurations of divine names: the wheat
manifests certain names but not others; the lion manifests another
configuration; the angel another. The human being, alone among
creatures, has the capacity to manifest all the divine names. The
human is the comprehensive locus.

This capacity is not automatic. The actual human falls short, in
varying degrees, of the comprehensive manifestation. The perfect
human is the rare individual whose locus has the comprehensive
manifestation actualized. The Prophet is the paradigm; the friends
of God (\arb{awliya'}) are subsequent partial manifestations.

\subsection{The cosmic role}

The doctrine of \arb{insan al-kamil} gives the human being a
cosmic role that has no parallel in Western philosophy of mind. The
perfect human is not just an intelligent organism; the perfect
human is the comprehensive locus of divine self-disclosure, the
mirror in which the divine sees itself reflected in its full
articulation.

This cosmic role grounds the Quranic doctrine of \arb{khalifah}
(vicegerent): the human is the divine representative on earth, the
one charged with manifesting the divine names in the structures of
created reality. The role is dignified and burdensome; the human
who fails in it is responsible for the failure.

For the economic project, the doctrine has implications we will
return to. The economic activity of the human is, on this view, the
manifestation of specific divine names (provision, justice,
generosity, exchange) in the structures of the material world. The
believer who organizes their economic life in accordance with the
divine names is partially actualizing the comprehensive
manifestation that the perfect human represents. The believer who
violates the divine structure (through riba, through theft, through
deception) frustrates the manifestation.

\subsection{The contemporary correspondence}

There is no direct correspondent in contemporary analytical
philosophy of mind to the doctrine of \arb{insan al-kamil}. The
contemporary frameworks have not articulated the cosmic role of
human consciousness in the way the Akbarian tradition has. This is
one of the places where the Akbarian framework adds substantial
content to the basic idealist position.

The closest contemporary parallel is perhaps the
``participatory anthropic principle'' developed by John Wheeler,
who argued that the universe requires observers to be fully real,
and that human consciousness has a constitutive role in the
universe.\footnote{Wheeler, J. A. (1990), ``Information, physics,
quantum: the search for links,'' in Zurek, W. H. (ed.),
\textit{Complexity, Entropy, and the Physics of Information},
Westview Press.} Wheeler's framework is much less worked out
than the Akbarian, but it is moving in a similar direction.

\section{Mulla Sadra: substantial motion}
\label{sec:mulla-sadra}

\subsection{Mulla Sadra's contribution}

Sadr al-Din Muhammad al-Shirazi (Mulla Sadra, d.~1641) was the
greatest Iranian philosopher of the Safavid period. His major work,
the \textit{al-Asfar al-arba`a al-`aqliyya} (The Four Intellectual
Journeys), is a systematic synthesis of Avicennan philosophy
(\arb{falsafa}), Suhrawardi's illuminationist philosophy (\arb{ishraq}),
and the Akbarian metaphysical tradition. Mulla Sadra was a Shi`i
Muslim and his influence has been particularly strong in Iranian
and Iraqi Shi`i intellectual life, but his synthesis is treated as
a major achievement across the broader Islamic philosophical
tradition.

The most distinctive Sadrian contribution is the doctrine of
\arb{harakat al-jawhariyyah}, ``substantial motion.''

\subsection{Substantial motion}

\begin{conceptbox}[\arb{Harakat al-jawhariyyah}]
\label{conc:substantial-motion}
The doctrine that all created beings are in constant ontological
flux. Things do not merely change accidentally (changing color,
position, age, etc.); they change substantially, in their very
being. The world is not a collection of fixed substances with
changing accidents; it is a continuous becoming.
\end{conceptbox}

The doctrine is a substantial development beyond the standard
Aristotelian conception, in which substantial change occurs only at
specific events (a substance ceases to be one kind and becomes
another, as when wood is burned and becomes ash). For Mulla Sadra,
substantial change is continuous: every created being is, at every
moment, in the process of becoming.

The doctrine integrates with the Akbarian \arb{tajalli}: the
continuous self-disclosure of \arb{wujud} produces, at every
moment, slightly different loci. The being is not the same from
moment to moment; it is being recreated. The continuity we
perceive is the continuity of the underlying process, not of any
substantial identity that persists across moments.

\subsection{Implications}

The doctrine of substantial motion has several implications for our
project.

First, it provides the metaphysical framework for understanding
\arb{tajalli} as not merely the maintenance of being but the
continuous transformation of being. The world is not just being
sustained; it is being continuously reconfigured. Each moment is
ontologically novel.

Second, it permits a richer account of dynamics. The standard
Aristotelian conception of substance and accident has difficulty
with becoming; substantial motion takes becoming as fundamental.
This is consistent with the dynamic conception of physics that
emerged in the twentieth century (relativity, quantum field theory)
and that has not been fully integrated with the static Aristotelian
ontology.

Third, it provides the metaphysical framework for the eschatological
content of Islamic tradition: the soul's progression toward its
ultimate state is not just accidental change but substantial
becoming. The believer who progresses spiritually is becoming a
different kind of being; the disbeliever who deteriorates is also
becoming a different kind of being. The eschatological judgment is
the recognition of what each soul has actually become.

\subsection{Integration with the Akbarian framework}

Mulla Sadra is not in tension with Ibn `Arabi; he is in continuity
with him. The Akbarian \arb{tajalli} provides the dynamic principle;
substantial motion provides the metaphysical analysis of what
happens to created beings under continuous \arb{tajalli}. The two
together constitute a coherent metaphysical system in which the
underlying unity of \arb{wujud}, the structural articulation through
divine names, the continuous manifestation through \arb{tajalli},
and the substantial becoming of creatures are all integrated.

This integrated system is what we will draw on in Part III. The
specific Islamic economic claims --- about \arb{rizq}, \arb{barakah},
\arb{niyyah} --- can be articulated within this system in ways
that the standard theological vocabulary alone cannot fully capture.

\section{Shah Wali Allah: economic theology}
\label{sec:shah-wali-allah}

\subsection{Shah Wali Allah's significance}

Shah Wali Allah Dehlawi (d.~1762) was the great eighteenth-century
Indian scholar. He was both a hadith scholar of the first rank ---
his commentary on the \textit{Muwatta} of Imam Malik is a classic
of the genre --- and a metaphysical thinker in the Akbarian
tradition. His major work, the \textit{Hujjat Allah al-Baligha}
(The Conclusive Argument from God), is a systematic treatment of
the rational basis of Islamic law and practice that draws extensively
on Akbarian ontology to articulate the metaphysical foundations of
legal and economic prescriptions.

For our project, Shah Wali Allah is the most directly relevant
classical thinker. He is the figure in the tradition who most
explicitly connects the metaphysical doctrines to the social and
economic order. His treatment of zakat, of inheritance, of riba,
and of waqf draws on the Akbarian framework to articulate
\emph{why} the Islamic prescriptions take the form they do.

\subsection{The structure of Hujjat Allah al-Baligha}

The \textit{Hujjat} is organized around the rational explication of
Islamic legal practices. For each major institution --- prayer,
fasting, zakat, pilgrimage, marriage, commerce --- Shah Wali Allah
provides a metaphysical and social analysis. The metaphysical
analysis grounds the institution in the structure of the divine
names and the dynamics of \arb{tajalli}. The social analysis
articulates the function the institution serves in the
\arb{ummah} (the community of believers).

The work is the closest classical Islamic source for the project of
this book. Where this book attempts to articulate the metaphysical
content of Islamic economic claims using contemporary frameworks,
Shah Wali Allah was attempting essentially the same project using
the Akbarian framework that was current in his time. The two
projects are continuous; the modern frameworks supplement, but do
not replace, the Akbarian articulation.

\subsection{Specific economic content}

We summarize briefly the relevant content from the \textit{Hujjat}.

\textbf{On zakat}: Shah Wali Allah articulates zakat as the
manifestation of the divine name \arb{ar-Rahman} (the Merciful) in
the social structure. The wealthy who pay zakat are not just
fulfilling a legal obligation; they are participating in the
manifestation of divine mercy in the social fabric. The poor who
receive zakat are receiving the divine mercy through the structure
that the wealthy have made themselves loci for. The economic
mechanism (redistribution) is the social manifestation of the
metaphysical reality (divine self-disclosure as merciful).

\textbf{On inheritance}: Shah Wali Allah analyzes the inheritance
shares as expressions of the divine names of justice (\arb{al-`Adl})
and wisdom (\arb{al-Hakim}). The specific shares (the wife receives
$1/8$ if there are children, $1/4$ if not; the male child receives
twice the female child's share; etc.) are not arbitrary; they
correspond to the structural realities of family responsibility
and dependency that the divine names articulate.

\textbf{On riba}: Shah Wali Allah's treatment is one of the most
sophisticated in the classical literature. He argues that riba
violates the structural principle by which exchange manifests the
divine names; it constitutes a structural injustice that destroys
the social fabric. The destruction is metaphysical, not just moral;
the social fabric that is destroyed is the loci of divine
self-disclosure in commerce.

\textbf{On waqf}: Shah Wali Allah treats waqf as the institution
that perpetuates the manifestation of divine generosity beyond the
lifetime of the donor. The waqf donor, in alienating property
permanently to a beneficiary purpose, is participating in the
divine name \arb{al-Karim} (the Generous) in a way that produces a
manifestation continuing across generations.

\subsection{The bridge to Part III}

These analyses by Shah Wali Allah are not the analyses we will give
in Part III, but they are the philosophical antecedents. Part III
will draw on the Akbarian framework articulated in this chapter and
on the contemporary frameworks of \xref{ch:quantum-foundations}
through \xref{ch:non-computability} to give updated, mathematically
precise treatments of the same Islamic economic phenomena that Shah
Wali Allah analyzed using the Akbarian framework alone. The
modernized treatments do not replace Shah Wali Allah's; they extend
them with the additional analytical resources now available.

\section{The contemporary correspondence}
\label{sec:contemporary-correspondence}

We now draw the explicit correspondence between the Akbarian
framework and the contemporary frameworks of Chapters 5 through 8.

\subsection{Wujud and Mind-at-Large}

The most direct correspondence is between absolute \arb{wujud} (the
foundational ontological category in Akbarian metaphysics) and
Mind-at-Large (the foundational ontological category in Kastrup's
analytical idealism). Both are single fundamental realities; both
underlie the apparent multiplicity of things; both are presupposed
by, rather than constructed from, the things that appear within
them.

The Akbarian framework adds the divine names (the structural
principles by which absolute \arb{wujud} differentiates) and
\arb{tajalli} (the active dynamic by which the differentiation
occurs). The analytical-idealist framework has the dissociation
analogy but not the structural articulation.

\subsection{Tajalli and quantum measurement}

A second correspondence: \arb{tajalli} as continuous self-disclosure
and quantum measurement as continuous actualization of the wave
function. Both are dynamic; both occur continuously; both are
required for the world to be as it appears.

The correspondence is structural; we are not claiming that quantum
measurement is \arb{tajalli}. We are claiming that the structural
account of ``continuous instantiation by an underlying process''
is articulated in both frameworks, and that the Akbarian framework
provides the metaphysical account of why the dynamic structure
described by quantum mechanics has the character it does.

\subsection{Wahdat al-wujud and the resolution of qadar}

A third correspondence: the Akbarian framework's account of how the
multiplicity of created things is consistent with the unity of
absolute \arb{wujud} provides the metaphysical structure within
which the Bohmian and QBist resolutions of \arb{qadar}
(\xref{ch:quantum-foundations}) are coherent. The unity of being
is the divine perspective; the multiplicity of created things is
the perspective of the dissociated alters; both are real; the
apparent tension is dissolved by the structural framework.

This is the fullest articulation. The contemporary frameworks
provide the formal apparatus; the Akbarian framework provides the
ontological substrate; together they constitute an account of the
classical \arb{qadar} problem that the classical \arb{kalam}
debates were attempting but did not have the resources to complete.

\subsection{The imaginal world and the limits of materialism}

A fourth correspondence: the Akbarian doctrine of the imaginal
world (\arb{`alam al-mithal}) and the contemporary recognition that
the materialist framework has trouble accommodating phenomena that
have aspects of both the mental and the physical (intentional
content, semantic meaning, phenomenal character of thought). The
Akbarian framework provides an ontological category for these
phenomena; the contemporary materialist framework does not.

The correspondence here is more suggestive than direct, but it
points to a place where the Akbarian framework adds substantial
content beyond what the contemporary analytical apparatus has
articulated. We expect this to be one of the productive directions
in subsequent volumes of the project.

\subsection{Insan al-kamil and the cosmic role of consciousness}

A fifth correspondence: the Akbarian doctrine of \arb{insan al-kamil}
(the cosmic role of human consciousness as comprehensive locus of
divine self-disclosure) and Wheeler's participatory anthropic
principle (the constitutive role of observers in the actualization
of physical reality). Wheeler's framework is much less developed
than the Akbarian, but it is moving in a similar direction.

Again, this is one of the places where the Akbarian framework adds
content beyond what contemporary philosophy has worked out. The
human being, on the Akbarian conception, has a cosmic role that the
contemporary materialist framework simply cannot articulate. The
contemporary participatory framework can articulate the structural
position; the Akbarian framework adds the theological and
existential content.

\section{Implications for the worked treatments of Part III}
\label{sec:akbarian-implications}

We end the chapter by previewing how the Akbarian framework will
appear in the worked treatments of Part III.

\subsection{Rizq as manifestation of ar-Razzaq}

Chapter 10 treats rizq as a conserved quantity. The Akbarian
framework adds: the conservation is not merely arithmetic; the
\arb{rizq} of each creature is the manifestation of the divine name
\arb{ar-Razzaq} in the locus of that creature. The integral of
\arb{rizq} over the creature's life is determined by the structural
character of the locus, which is fixed at the moment of creation
(\xref{sec:wujud-foundational}). The variations within the life are
the temporal distribution of the manifestation; the integral is the
total structural capacity of the locus for receiving the divine
self-disclosure.

\subsection{Children-bring-rizq as expanded locus}

Chapter 11 treats the network-theoretic account of the
``children-bring-rizq'' claim. The Akbarian framework adds: a
larger family is a larger locus for the manifestation of
\arb{ar-Razzaq}. The added members do not divide a fixed
manifestation; they expand the capacity of the locus to receive
the manifestation. The super-linear scaling of the formal account
is the formal expression of this metaphysical structure.

\subsection{Halal vs.\ haram as structural alignment}

Chapter 12 treats the ergodic vs.\ non-ergodic distinction between
halal and haram earnings. The Akbarian framework adds: halal
earnings are configurations of exchange that are structurally
aligned with the divine names; haram earnings are misaligned. The
formal ergodicity (or non-ergodicity) is the structural consequence
of this alignment (or misalignment).

\subsection{Riba as structural destruction}

Chapter 13 treats riba as financial entropy. The Akbarian framework
adds: riba is a structural violation of the divine names that
articulate just exchange. The structural violation is what produces
the destruction of real wealth that the formal analysis identifies.

\subsection{Gratitude as alignment of locus}

Chapter 14 treats gratitude as the multiplier of wealth. The
Akbarian framework adds: gratitude is the disposition that aligns
the human locus with the source of \arb{tajalli}. The aligned
locus receives the manifestation more fully; the misaligned locus
receives it less fully. The behavioral mediation chain identified
in the formal analysis is the operational mechanism of this
alignment.

\subsection{Barakah as intensity of manifestation}

Chapter 15 treats barakah as the operational multiplier on material
wealth. The Akbarian framework adds: barakah is the intensity with
which the divine name (typically \arb{ar-Razzaq} but also others)
is manifested in the locus of the wealth. Two material configurations
of the same nominal magnitude can have different barakah because
they are different intensities of manifestation.

\subsection{Qadar as the unity of being}

Chapter 16 treats qadar with the Bohmian and QBist formalisms. The
Akbarian framework adds: the divine knowledge that is perfect is
the knowledge of \arb{wujud} of itself, of which all created
realities are loci. The human knowledge that is finite is the
knowledge of one dissociated alter. Both are real; the apparent
tension is dissolved by the structural framework.

\subsection{Anti-concentration as preservation of structure}

Chapter 17 treats the four anti-concentration mechanisms (zakat,
inheritance partition, riba prohibition, waqf alienation) as a
coordinated package. The Akbarian framework adds: these mechanisms
together preserve the social structure as a manifestation of the
divine names. Concentration of wealth in a few hands corresponds to
the manifestation collapsing onto narrow loci; the preservation
mechanisms keep the manifestation distributed across the
articulated structure.

\subsection{The unified picture}

Across all eight worked treatments, the Akbarian framework provides
the metaphysical substrate. The formal frameworks (network theory,
ergodicity economics, behavioral mediation, etc.) provide the
operational mechanisms. The two together constitute treatments that
are both rigorously formal and metaphysically deep.

This is what we mean by Level 3 treatment in the strongest sense
(\xref{ch:levels-of-rigor}): the metaphysical claim corresponds to
a real mechanism, the mechanism is mathematically derivable, and
the metaphysical content is not collapsed into the mechanism but
informs and enriches the operational analysis.

\section{Conclusion of Chapter 9 and of Part II}

Chapter 9 has developed the Akbarian metaphysical tradition at the
depth required for the rest of the project. The reader has been
shown:

\begin{itemize}[leftmargin=2em]
    \item The Akbarian conception of \arb{wujud} as the foundational
    ontological category.
    \item The doctrine of \arb{wahdat al-wujud}, carefully
    distinguished from pantheism.
    \item The divine names as principles of articulation.
    \item \arb{Tajalli} as the active dynamic of continuous creation.
    \item The imaginal world as intermediate ontological domain.
    \item The perfect human as comprehensive locus.
    \item Mulla Sadra's substantial motion.
    \item Shah Wali Allah's economic-theological synthesis.
    \item The correspondences with the contemporary frameworks of
    Chapters 5 through 8.
    \item The previews of how the Akbarian framework will appear in
    the worked treatments of Part III.
\end{itemize}

Part II is now complete. The reader has been given the modern
analytical apparatus (Chapters 5--8) and the Akbarian metaphysical
substrate (Chapter 9) that the worked treatments of Part III will
require. The integration of the two --- formal analysis grounded in
metaphysical depth --- is what makes the project distinctive.

We are now prepared for Part III. Each of the eight chapters there
will take one canonical Islamic metaphysical claim, articulate it
precisely, formalize it using the apparatus of Part II, derive its
content, and report the empirical and theoretical implications. The
treatments will be the most substantive in the book. The reader who
has worked through Part II is now equipped to follow them.

\begin{summarybox}[Chapter summary]
\textbf{Ibn al-`Arabi}: the central figure of the Akbarian tradition;
his \textit{Futuhat al-Makkiyyah} and \textit{Fusus al-Hikam} are
the foundational texts; his philosophical influence has been
substantial across the Muslim world for eight centuries.

\textbf{\arb{Wujud}}: the foundational ontological category; admits
of degrees; absolute \arb{wujud} is divine, the source of all
derived \arb{wujud}.

\textbf{\arb{Wahdat al-wujud}}: the doctrine that there is only one
\arb{wujud}; created things are loci of divine self-disclosure;
distinguished carefully from pantheism.

\textbf{The divine names}: structural principles by which absolute
\arb{wujud} unfolds into determinate beings; each created thing is
the locus of a configuration of divine names.

\textbf{\arb{Tajalli}}: theophany, self-disclosure; the active
dynamic by which absolute \arb{wujud} manifests in finite loci;
continuous, not once-and-for-all.

\textbf{The imaginal world}: intermediate ontological domain
between pure spirit and pure matter; the medium through which
spirit and matter interact; accessed through imagination as a real
cognitive faculty.

\textbf{\arb{Insan al-kamil}}: the perfect human; the comprehensive
locus that manifests all the divine names; the cosmic role of
human consciousness.

\textbf{Mulla Sadra}: substantial motion; all created beings are in
constant ontological flux; integrates with the Akbarian
\arb{tajalli}.

\textbf{Shah Wali Allah}: the most explicit classical articulation
of the connection between Akbarian metaphysics and Islamic economic
prescriptions; the closest classical antecedent to the present
project.

\textbf{Contemporary correspondences}: \arb{wujud} and Mind-at-Large;
\arb{tajalli} and quantum measurement; \arb{wahdat al-wujud} and the
resolution of \arb{qadar}; the imaginal world and the limits of
materialism; \arb{insan al-kamil} and Wheeler's participatory
anthropic principle.

\textbf{For Part III}: the Akbarian framework provides the
metaphysical substrate that the formal treatments of the eight
canonical claims will operate on. The integration of formal
analysis and metaphysical depth is what the project is attempting.

\textbf{What comes next}: Part III. Eight chapters, each treating
one canonical Islamic metaphysical claim with the apparatus of
Parts I and II.
\end{summarybox}

\cleardoublepage
